\section{Conclusions}

The horizon gate and the cross-covariance penalty split the persistent factor
into a designated coordinate block, and across all six settings formed by three
datasets and two loss families they jointly produce a separation that neither
device reaches alone. The separation requires a latent target: switching only
the prediction target to the raw signal halves the separation score, and what
breaks down is precisely the inclusion of the persistent factor. On synthetic
data we exceed label-free post-hoc rotation for both loss families, but on real
data we fall short of Slow Feature Analysis, losing three of the four real
settings, and the separation costs $0.111$ macro-F1 downstream on HAPT. The
gate trades absolute performance for stability under domain shift.

\subsection{Relevance to Responsible AI}

\textbf{Explanation comes from the design, not after the fact.} Most
explainability techniques attach a story---saliency maps, attention
visualizations---on top of an already trained entangled model. We instead put
the axis of interpretation inside the training objective, so $\zslow$ is
reserved for the persistent state before any probe is run, as measured by the
block $\times$ factor matrix and by SEP. But the boundary is plain:
the block structure is interpretable, the encoder is not. The latent space is
rotation symmetric, so any 16 dimensions can hold the persistent factor; we
assign no meaning to individual dimensions.

\textbf{It degrades less where it matters.} On HAPT, split by subject,
$\zslow$ drops the least when the domain is perturbed and labels are
scarce---at 25 labels it loses $0.113$ less macro-F1 than a width-matched
mechanism-free block, five times the seed standard deviation. But that is
degrading less, not being better: absolute performance never overtakes the
mechanism-free embedding at any perturbation strength tested, and the effect
exceeds the error bars in only one of six settings.

\textbf{The cost becomes a measurable quantity.} \emph{The design makes the
downstream cost measurable}: because the method explicitly produces a block
holding the persistent state, the cost of solving a task from that block alone
is a single number---the macro-F1 difference against the full embedding
$z_{\mathrm{full}}$ of a model trained with no mechanism at all. With an
entangled representation there is no such quantity to measure.

\subsection{Limitations}

\textbf{First, the hyperparameters cannot be chosen without labels.} The
cross-covariance coefficient, the block width, the target width, the patch
size, and the horizon range have no derivation rule and were set by downstream
performance. The autocorrelation-derived horizon threshold is a partial
exception, but it fails on regularly repeating signals: on PTB-XL the per-record estimates scatter from $0.1$ to $49.9$ patches and the rule returns $56.6$, outside the training horizon range, so we set the threshold by hand instead. The exclusion term needs no task labels, so a failed separation can at
least be detected label-free.

\textbf{Second, we did not analyze the target width.} We used 8, 12, and 8
patches, but the ratio to the transient factor's autocorrelation time is
$1.53$, $0.91$, and $0.42$---not constant. This width governs how much of the
transient factor is averaged away inside the target; the sweep was not run.

\textbf{Third, we did not demonstrate a practical benefit:} the perturbed-domain effect clears the error bars in only one of six settings, and even there absolute performance never overtakes the mechanism-free embedding.

\textbf{Fourth,} the claim that a factor can be ``persistent yet absent from
the slow part of the signal'' rests on a single synthetic dataset. One
counterexample refutes a universal proposition, but we designed it ourselves,
and neither PTB-XL nor HAPT exhibits the case.

\textbf{Fifth, our axis selects persistence, not relevance.} A persistent
nuisance factor---the low-frequency drift respiration and electrode motion
produce in an electrocardiogram---enters $\zslow$ by the same criterion, and
which persistent factors matter remains a human judgment.

\subsection{Future Work}

We compared only against label-free post-hoc rotation, so methods that place factor separation explicitly in the objective remain to be compared. Validation is also needed on data whose persistent factor varies within a recording while staying constant within a window --- PTB-XL fails the first property and HAPT the second, whereas sleep-stage classification has both. Finally, using the trajectory of the separated block as a state indicator, rather than as classifier input, would address the fourth limitation directly.
